% Глава: 07-text-pipelines
% Статус: наполнено (первый прогон).
\chapter{Стандартные потоки, перенаправление, конвейеры}
\label{ch:07-text-pipelines}

\section{Стандартные потоки}

Каждый запускаемый процесс по умолчанию получает три
\textbf{стандартных потока}:

\begin{center}
\begin{tabular}{clp{0.4\linewidth}}
\toprule
\textbf{FD} & \textbf{Поток} & \textbf{Назначение} \\
\midrule
0 & stdin  & вход: то, что программа читает \\
1 & stdout & выход: результат работы программы \\
2 & stderr & ошибки и диагностика \\
\bottomrule
\end{tabular}
\end{center}

По умолчанию все три подключены к терминалу. Ключевой принцип
Unix: программы пишут результат в \cmd{stdout}, а сообщения
об ошибках --- в \cmd{stderr}. Разделение позволяет
независимо перенаправлять данные и диагностику.

\section{Перенаправление вывода}

\begin{minted}[bgcolor=gray!6]{bash}
ls > out.txt        # stdout в файл: создаёт или очищает
ls >> out.txt       # stdout добавляется в файл
ls 2> err.txt       # stderr в файл
ls &> all.txt       # stdout и stderr в один файл (bash)
ls 2>&1 > out.txt   # stderr туда, где сейчас stdout, потом stdout в файл
ls > out.txt 2> err.txt    # явно разделить
tee out.txt < big.txt       # показать на экране и в файле
\end{minted}

Запомнить два оператора: \cmd{>} --- заменить/создать файл,
\cmd{>>} --- дописать.

\begin{warnbox}
\cmd{>} \textbf{обрезает} файл до нуля перед записью. Если
файл уже содержал важные данные и вы хотели дописать,
используйте \cmd{>>}.
\end{warnbox}

\section{Конвейеры (pipes)}

Оператор \cmd{|} передаёт \cmd{stdout} одной команды в
\cmd{stdin} другой. Это основной способ строить сложные
преобразования из простых:
\begin{minted}[bgcolor=gray!6]{bash}
cat access.log | grep "GET" | wc -l
df -h | grep " /$"
ps aux | sort -k3 -nr | head
\end{minted}

Правило: каждое звено конвейера делает \textbf{одну} операцию.
Если решение требует трёх отдельных команд с временными
файлами --- скорее всего, его можно выразить конвейером.

\section{Текстовые утилиты}

\begin{center}
\begin{tabular}{lp{0.6\linewidth}}
\toprule
\textbf{Утилита} & \textbf{Назначение} \\
\midrule
\cmd{cat}   & вывести файл(ы) в stdout; склеить файлы \\
\cmd{head}  & первые строки/байты (\flag{n} 5 --- 5 строк) \\
\cmd{tail}  & последние строки; \flag{f} --- следить за обновлением \\
\cmd{wc}    & подсчитать строки/слова/байты (\flag{l}, \flag{w}, \flag{c}) \\
\cmd{sort}  & отсортировать строки (\flag{r} --- обратно, \flag{n} --- число, \flag{k} --- ключ) \\
\cmd{uniq}  & убрать подряд идущие дубли (работает после sort) \\
\cmd{cut}   & вырезать поле по разделителю (\flag{-d} '\textbackslash|' \flag{-f} 2) \\
\cmd{grep}  & поиск по регулярному выражению \\
\cmd{tr}    & побайтовая замена/удаление символов \\
\cmd{less}  & постраничный просмотр больших файлов \\
\bottomrule
\end{tabular}
\end{center}

\cmd{grep} --- самая частая. Основные флаги:
\begin{minted}[bgcolor=gray!6]{bash}
grep ERROR log.txt         # строки, содержащие ERROR
grep -i error log.txt      # без учёта регистра
grep -c ERROR log.txt      # сколько таких строк
grep -E "err|warn" log.txt # расширенные регулярные выражения
grep -v DEBUG log.txt      # строки, НЕ содержащие DEBUG
grep -n "user" config      # с номерами строк
\end{minted}

\section{Типовой сценарий анализа}

\begin{sloppypar}
На практике такой лог-файл обычно лежит на вашей ВМ на сервере
\fpath{pglab.nsu.ru}. Допустим, в \fpath{events.log} строки вида
\cmd{timestamp|level|source|file|status}.
\end{sloppypar}
Частые вопросы:
\begin{minted}[bgcolor=gray!6]{bash}
wc -l < events.log                       # общее число строк
grep -c ERROR events.log                 # число ошибок
cut -d'|' -f2 events.log | sort | uniq -c # количество каждого уровня
cut -d'|' -f3 events.log | sort | uniq -c | sort -nr # по источникам
grep ERROR events.log > errors.log        # выписать ошибки
\end{minted}

Ключевой приём: \cmd{sort} перед \cmd{uniq} ---
\cmd{uniq} убирает только \emph{соседние} дубли, поэтому
сначала строки должны быть отсортированы.

\section{Постраничный просмотр: \cmd{less}}

Для больших файлов \cmd{cat} неудобен: вывод проматывает
терминал. \cmd{less} --- стандартный pager:
\begin{minted}[bgcolor=gray!6]{bash}
less big.log
# внутри less:
# /pattern  n  N  - поиск вперёд/назад
# g  G            - начало / конец
# Space / b       - страница вперёд / назад
# q               - выход
\end{minted}

\cmd{less} и \cmd{more} похожи, но \cmd{less} позволяет
прокручивать назад и искать.

\section{Типичные ошибки}
\begin{itemize}
  \item Конвейер «ничего не возвращает», хотя отдельные
    команды работают --- проверьте, не фильтрует ли
    \cmd{grep} всё (нет совпадений);
  \item \cmd{uniq} не убирает все дубли --- не отсортированы
    строки перед \cmd{uniq};
  \item \cmd{>} случайно перезаписал важный файл;
  \item Забыли \cmd{2>} и диагностика ушла в терминал,
    а не в файл;
  \item \cmd{tail} без \flag{f} показывает конец один раз
    и завершается; для наблюдения нужен \flag{f}.
\end{itemize}

\section{Связь с остальным курсом}

Конвейеры --- это основа автоматизации: из них собираются
скрипты оболочки (\nameref{ch:27-scripting-ide}),
обрабатываются логи приложений, выводятся результаты
тестов (\nameref{ch:15-testing-optimization}), логируются
запуски контейнеров (\nameref{ch:19-docker-basics}).

\section{Задачи для самопроверки}
\begin{enumerate}
  \item Объясните разницу между \cmd{>} и \cmd{>>}. Что
    произойдёт, если вы случайно использовали первое?
  \item Что такое \cmd{stderr} и зачем его разделять от
    \cmd{stdout}?
  \item Напишите конвейер, который считает количество
    строк в файле, содержащих слово \fpath{error}
    (без учёта регистра).
  \item Почему \cmd{sort} обычно ставят перед \cmd{uniq}?
  \item Как посмотреть последние 20 строк журнала и
    продолжать наблюдать за новыми?
\end{enumerate}
